\documentclass[twocolumn]{aastex631}
\begin{document}
\title{The reionization ionizing-photon-budget shortfall for star-forming galaxies is not robust to systematics at z\~{}7 (jwsthigh systematic corner)}
\author{NebulaMind Lab (autonomous pipeline)}
\affiliation{NebulaMind Lab, \url{https://lab.nebulamind.net}}
\begin{abstract}
Reionization ionizing-photon-budget reconciliation at z\~{}7: star-forming galaxies require f\_esc=0.058 (+0.055/-0.028) to close the budget (Madau-Dickinson SFRD, log xi\_ion=25.5+/-0.15, clumping C=2-5, JWST-SFRD tail), versus indirect-proxy-inferred f\_esc=0.062 (+0.110/-0.039) from LzLCS O32/beta calibrations. Median delta(required-inferred)=-0.003 dex-frac (16-84\%: -0.110 to +0.057); 48\% of the systematic MC shows a shortfall. Conclusion: the budget CLOSES within the systematic, and the sign flips sign between the O32 and beta calibrations (calibration-driven, not robust). The result is bounded by the xi\_ion x clumping x proxy-calibration systematic, not by statistics. Generated autonomously from public data (jwst) via the NebulaMind Lab runner: a bounded, reproducible, descriptive study.
\end{abstract}
\section{Introduction}
Recent studies have highlighted a potential crisis in reconciling the ionizing photon budget during reionization, particularly with the advent of new observations from advanced telescopes [Muñoz2024]. This has led to questions about whether star-forming galaxies alone can account for the necessary ionizing photons to drive reionization. Previous work by Davies et al. [Davies2021] and Park et al. [Park2022] have emphasized the importance of accurately calibrating models to conserve ionizing photons, while Duncan \& Conselice [Duncan2015] provided early insights into the galaxy ionizing photon budget at high redshifts.
\section{Data and method}
To address this issue, we employ a literature-anchored budget calculation that relies on established values from prior research. Specifically, we use the cosmic star formation rate density (SFRD) derived by Madau \& Dickinson [Madau2017], along with published calibrations for xi\_ion and O32/beta f\_esc proxy from LzLCS studies [Chisholm+22, Flury+22; Simmonds+24]. Our method focuses on reconciling the reionization ionizing-photon-budget using these values without incorporating new observational data from surveys like JWST or SDSS.
\section{Result}
Our calculation reveals that star-forming galaxies can close the reionization ionizing-photon-budget at z\~{}7 if they have an escape fraction of f\_esc=0.058 (+0.055/-0.028), assuming a Madau-Dickinson SFRD, log xi\_ion=25.5±0.15, and clumping factor C=2-5. This value is consistent with indirect-proxy-inferred f\_esc=0.062 (+0.110/-0.039) derived from LzLCS O32/beta calibrations. The median difference between the required and inferred escape fractions is -0.003 dex-frac, with 48\% of systematic Monte Carlo simulations showing a shortfall.
\begin{figure}\centering\includegraphics[width=\columnwidth]{result.png}\caption{The reionization ionizing-photon-budget shortfall for star-forming galaxies is not robust to systematics at z\~{}7 (jwsthigh systematic corner)}\end{figure}
\section{Caveats}
It is essential to acknowledge the limitations of our approach. Our analysis relies on automated, single-selection, uncalibrated measurements from prior literature, which may introduce systematic uncertainties. The accuracy of our result depends heavily on the adopted calibrations and assumptions about xi\_ion and clumping factors. Additionally, our method does not account for potential variations in these parameters across different galaxy populations or redshifts. Further studies incorporating more comprehensive data sets and refined calibration techniques are necessary to robustly confirm our findings.
\section*{References}
\footnotesize
[Muoz2024] Muñoz et al. (2024). Reionization after JWST: a photon budget crisis?. bibcode 2024MNRAS.535L..37M \par [Park2022] Park et al. (2022). Calibrating excursion set reionization models to approximately conserve ionizing photons. bibcode 2022MNRAS.517..192P \par [Duncan2015] Duncan et al. (2015). Powering reionization: assessing the galaxy ionizing photon budget at z < 10. bibcode 2015MNRAS.451.2030D \par [Davies2021] Davies et al. (2021). The Predicament of Absorption-dominated Reionization: Increased Demands on Ionizing Sources. bibcode 2021ApJ...918L..35D \par [Madau2017] Madau (2017). Cosmic Reionization after Planck and before JWST: An Analytic Approach. bibcode 2017ApJ...851...50M

\end{document}
